% SPDX-License-Identifier: LPPL-1.3c
% Copyright (C) 2026 Christophe Jorssen
% Author and Current Maintainer: Christophe Jorssen <christophe.jorssen@gmail.com>
% LPPL maintenance status: maintained
% This file is part of luafemm. See LICENSE and LICENSES.md for its terms.
\section{Refining a mesh from the computed field}
\label{sec:adaptive-mesh}

Start with a coarse mesh, complete the geometry, then request adaptation before
capturing final profiles or measurements. The estimator combines element current
residuals and jumps of tangential H across shared edges. Bulk marking selects
triangles accounting for a specified fraction of the total squared indicator.
The present indicator includes interior jumps and volume currents, not a separate
weak-boundary residual; it is heuristic for nonlinear and anisotropic media.
Marked triangles receive red refinement; their neighbours receive conforming
green splits. Material interfaces, holes and narrow features keep their original
coordinates. This procedure cannot repair a poorly approximated curved boundary.

\begin{command}{\femmadapt\oarg{options}}
Solve and locally refine the current planar model. Ordinary static and harmonic
P1 models are supported; imported states, periodicity, operating responses,
motion, external networks and flux controllers are excluded. Use |cache mode=off|
(the default); save a finished answer as |.ansx| if persistence is needed.
Each successful pass replaces the current model in place. A failed next pass
leaves the last solved mesh available. Previously captured profiles remain owned
snapshots and are not retroactively updated.
\end{command}
\begin{key}{/tikz/femm/adapt=\meta{options}}
Invoke |\femmadapt| through a TikZ key, after all physical path postactions.
Do not put this operation among the options of a region still being constructed.
\end{key}
\begin{key}{/femm/adapt/passes=\meta{integer} (initially 3)}
Maximum additional refinements, from 1 to 12. The initial solve is separate.
Reaching this limit returns the last field and indicator, not a claim of accuracy.
\end{key}
\begin{key}{/femm/adapt/fraction=\meta{number} (initially .5)}
Fraction in $(0,1]$ of the summed squared indicators to mark per pass.
Larger values spread refinement more broadly; 1 marks the entire positive set.
\end{key}
\begin{key}{/femm/adapt/tolerance=\meta{nonnegative number} (initially .05)}
Stop when the dimensionless residual indicator divided by the field-energy norm
falls below this value. Zero requests the stated passes unless the residual is
exactly zero. This is an estimator, \emph{not} a certified error percentage.
\end{key}
\begin{key}{/femm/adapt/maximum nodes=\meta{integer} (initially 100000)}
Hard node budget, between 3 and 1000000. Budget and minimum-angle failures occur
before replacing the previous solved model. The existing |minimum angle| remains
an acceptance condition; reduce it explicitly if it conflicts with local splits.
\end{key}

Use |\femmstat| with the following names:
\begin{center}
\begin{tabular}{ll}
\toprule Statistic & Meaning\\
\midrule
|adaptive_error| & Final normalized residual indicator\\
|adaptive_passes| & Completed additional refinements\\
|adaptive_seconds| & CPU time in seconds\\
|adaptive_memory_kib| & Nonnegative Lua heap increase in KiB\\
\bottomrule
\end{tabular}
\end{center}
The heap increase is not peak process memory.
The mesh statistics |min_angle| and |max_area| describe the final mesh.
Localized problems may benefit substantially; a smooth global field may favour
uniform refinement. Convergence of flux, loss and especially force must still be
checked directly, with geometric approximation and exterior extent held separately.

\begin{codeexample}[width=5cm]
\begin{tikzpicture}[femm/problem={xmin=-20,xmax=20,
 ymin=-20,ymax=20,mesh size=5}]
 \path[femm/region={material=air,current density=1e6}]
   (-6,-6) rectangle (6,6);
 \femmadapt[passes=2,fraction=.6,tolerance=0,
   maximum nodes=5000]
 \pic {femm scalar map};
 \pic {femm mesh};
\end{tikzpicture}
\end{codeexample}
